\subsection{QEMU Object Model Integration}
\label{sec:imp_qom_integration}

The STM32F4xx I$^2$C controller emulation is integrated into QEMU's Object Model (QOM) through the following declarations:

\begin{lstlisting}[
    language=C,
    caption={QOM macros},
    label={code:qom_macros_1}]
#define TYPE_STM32F4XX_I2C "stm32f4xx-i2c"
OBJECT_DECLARE_SIMPLE_TYPE(STM32F4XXI2CState, STM32F4XX_I2C)
\end{lstlisting}

\textbf{\texttt{TYPE\_STM32F4XX\_I2C}} defines the unique QOM type identifier string 
\\ \texttt{"stm32f4xx-i2c"}, enabling device instantiation via \texttt{qdev\_new(TYPE\_\\STM32F4XX\_I2C)} and \texttt{object\_initialize\_child()} calls throughout the SoC emulation.
\textbf{\texttt{OBJECT\_DECLARE\_SIMPLE\_TYPE}} establishes the type-safe relationship between the \texttt{STM32F4XXI2CState} instance structure and its QOM type, automatically generating casting macros.

\subsubsection{Device Lifecycle}
\label{subsec:lifecycle}

The first thing QEMU does is register the device within its kernel. This is done with the \texttt{type\_init} macro. This macro is called before the main Qemu macro.

\begin{lstlisting}[
    language=C,
    caption={\texttt{type\_init} macro},
    label={code:type_init}]
type_init(stm32f4xx_i2c_register_types)
\end{lstlisting}

The \texttt{TypeInfo} structure that represents the I2C controller is the following

\begin{lstlisting}[
    language=C,
    caption={\texttt{TypeInfo} structure},
    label={code:TypeInfo}]
static const TypeInfo stm32f4xx_i2c_info = 
{
    .name           = TYPE_STM32F4XX_I2C,
    .parent         = TYPE_SYS_BUS_DEVICE,
    .instance_size  = sizeof(STM32F4XXI2CState),
    .class_init     = stm32f4xx_i2c_class_init,
};
\end{lstlisting}

Upon device registration, QEMU invokes the \textbf{\texttt{stm32f4xx\_i2c\_init()}} function, Code~\ref{code:stm32f4xx_i2c_class_init} . 
This class initializer overrides the virtual methods of \texttt{SysBusDeviceClass} with 
STM32F4xx I2C-specific implementations, configures device properties via \texttt{device} \texttt{\_class\_set\_props()}, assigns the realization callback \texttt{stm32f4xx\_i2c\_realize}, 
and registers the VM state description. In this specific implementation, the \textbf{\texttt{stm32f4xx}} \textbf{\texttt{\_i2c\_init()}} function it also allocates the FSM structure.
    
\begin{lstlisting}[
    language=C,
    caption={\texttt{stm32f4xx\_i2c\_init()}},
    label={code:stm32f4xx_i2c_class_init}]
static void stm32f4xx_i2c_class_init(ObjectClass *klass, void *data)
{
    DeviceClass *dc = DEVICE_CLASS(klass);

    device_class_set_props(dc, stm32f4xx_i2c_properties);
    dc->realize = stm32f4xx_i2c_realize;
    dc->vmsd    = &vmstate_stm32f4xx_i2c;
}
\end{lstlisting}

